\section{Related work}
\label{sec:related}

\paragraph{Multi-agent benchmarks.}
Pommerman offers a two-versus-two grid game with a radio variant that allows discrete messages between teammates \citep{resnick2018pommerman}. Territory Paint Wars uses simultaneous tile competition between two players to diagnose failures of PPO self-play \citep{singh2026paintwars}. Melting Pot 2.0 evaluates policies alongside unfamiliar co-players across environments that include Territory \citep{agapiou2022meltingpot}. MARL-BattleGrounds pairs team matches with short tactical endgame tests on held-out positions \citep{hawili2026battlegrounds}. Territory Control keeps the transition rules fixed across protocols and varies whether agents can observe teammates' committed actions or tentative proposals for the current tick. Replays let us measure the effect of that information on a common set of tactical opportunities.

\paragraph{Sequential decision making.}
\citet{ding2024seqcomm} propose SeqComm, which assigns decision priority by exchanging encoded observations and evaluating predicted intentions. Agents then pass committed actions to teammates that decide later, and the environment executes the resulting actions simultaneously. Our protocol fixes the within-team order by slot and exposes the prefix of committed actions on a frozen board. This makes it possible to mask the prefix while holding decision order fixed. The Multi-Agent Transformer also generates actions in sequence, using one encoder-decoder over the agents' observations and an autoregressive joint-action decoder \citep{wen2022mat}. In our protocol each actor call takes its own observation and the commits the environment has made visible to that slot.

\paragraph{Intention sharing and message gradients.}
\citet{kim2021intention} use learned dynamics and models of teammates to generate and communicate compressed future trajectories. In our two-round protocol, each agent sends one of nine actions for the \emph{current} tick, receives the team's first-round proposals, and can revise its own action before execution. Multi-task Communication Skills (MCS) couples a learned message representation to the sender's executed actions through a prediction objective \citep{zhu2026mcs}. We use this objective to predict a sender's second-round action from its first-round action index. DIAL backpropagates through a real-valued channel during centralized training and quantizes messages at execution \citep{foerster2016dial}. Our receiver sees the recorded discrete proposal in the forward pass. A straight-through estimator lets us vary whether its policy gradient reaches the sender's proposal distribution while retaining the same proposal format.

\paragraph{Evaluating communication and cooperation.}
SMACv2 introduced procedurally generated scenarios after showing that policies conditioned only on time could win some SMAC games \citep{ellis2023smacv2}. Coordination diagnostics can also reveal differences in interference and joint task completion between policies with similar returns \citep{cardei2026coordination}. \citet{lowe2019pitfalls} distinguish \emph{positive signaling}, in which a message carries information about a sender, from \emph{positive listening}, in which that message changes a receiver's behavior. We test listening with a fixed checkpoint and first-round proposals and count the consequences for cover and movement. Full matches capture how these decisions affect the positions each policy goes on to visit.
